\chapter{PTMCMC: Parallel Tempering}
\label{ch:ptmcmc}

\newthought{This chapter assumes Chapter~\ref{ch:toymcmc}.} If you have not read it,
do that first --- the Metropolis--Hastings step, the class hierarchy, deterministic
\textsc{uuid}s, checkpointing, and the \textsc{mcmc}-family outputs are all identical
here. This chapter covers exactly one thing: what changes when several Metropolis
chains run at different temperatures and periodically swap, and how to configure and
read that.

\section{What parallel tempering adds}
\label{sec:pt-what}

A single random-walk chain (Chapter~\ref{ch:toymcmc}) can get stuck: if the posterior
has two modes separated by a region of very low likelihood, the chain essentially
never proposes a step through that region, and a chain that started in one mode may
never discover the other.

Parallel tempering\cite{EarlDeem:2005} runs \yamlkey{num\_chains} replicas of the
\emph{same} chain in parallel, each at its own temperature $T_i$ (inverse temperature
$\beta_i = 1/T_i$), and periodically lets adjacent-temperature replicas swap their
current state:

\begin{itemize}
  \item The \textbf{cold} replica, $T=1$, sees the real posterior --- this is the one
        whose points you actually want.
  \item \textbf{Hot} replicas, $T>1$, see a flattened likelihood surface
        ($L^{\beta}$ with $\beta<1$): low-likelihood barriers that trap the cold chain
        are far easier to cross at high temperature.
  \item An accepted \textbf{exchange} hands the hot replica's newly discovered region
        one rung down the ladder, eventually reaching the cold replica --- without the
        cold chain ever having to propose the crossing step itself.
\end{itemize}

\begin{jwarn}
\textbf{Only the cold replica is your posterior.} \Jtwo\ archives every replica's
every step, at every temperature (\S\ref{sec:ptmcmc-outputs}) --- but the hot replicas'
points are draws from a deliberately flattened surface, not the posterior you asked
for. Chain \textsc{ids} are assigned in ladder order and
\yamlkey{temperature\_ladder[0]} is always the required cold value $1.0$
(\S\ref{sec:ptmcmc-yaml}), so \textbf{\yamlkey{chain\_id} 0 is always the cold,
posterior-valid replica} --- filter \file{samples.csv} to \yamlkey{chain\_id} \code{==
0} before treating it as a posterior sample.
\end{jwarn}

\section{The exchange move}
\label{sec:pt-algorithm}

Every \yamlkey{exchange\_interval} draws, adjacent rungs on the ladder attempt a swap:

\begin{enumerate}
  \item \textbf{Pair adjacent rungs.} Only neighbouring temperatures are ever paired
        --- $(T_0,T_1), (T_2,T_3), \ldots$ on one round, offset by one rung on the
        next, so every adjacent pair gets a turn over time. The hottest and coldest
        replicas are never paired directly.
  \item \textbf{Accept or reject the swap.} With $\beta_i=1/T_i$, $\beta_j=1/T_j$ and
        current log-likelihoods $\log L_i$, $\log L_j$, let
        $\Delta = (\beta_i-\beta_j)(\log L_j-\log L_i)$. Accept (swap) if
        $\Delta \geq 0$, or with probability $e^{\Delta}$ otherwise --- the same
        Metropolis shape as an ordinary step, applied between chains instead of within
        one.
  \item \textbf{Swap state, not temperature.} An accepted exchange swaps the two
        replicas' parameter vectors and log-likelihoods. \textbf{The temperature stays
        with the ladder slot}, not with the configuration --- rung $i$ is always at
        $T_i$; what moves between rungs is which configuration currently occupies
        each one.
\end{enumerate}

\begin{jnote}
This is the detail that trips people up: parallel tempering does \emph{not} mean
"each chain's temperature changes over time." The ladder is fixed; configurations
migrate along it. \yamlkey{chain\_id} \code{0} is permanently the cold rung for the
whole run --- it never itself "becomes" a hot replica.
\end{jnote}

\section{Why the swap rule works: detailed balance on the extended ensemble}
\label{sec:pt-detailed-balance}

Chapter~\ref{ch:toymcmc}'s derivation (\S\ref{sec:mcmc-detailed-balance}) already
covers why an ordinary Metropolis step leaves a target distribution invariant; the
only change here is which target each replica's step leaves invariant. Replica $i$
runs exactly that step against $\pi_i(u) \propto L(u)^{\beta_i}$ --- the general
$\beta$ that step~3 there fixed at $1$ for \sampler{ToyMCMC} is now the rung's own
$\beta_i$. Because $\log \pi_i(u) = \beta_i \log L(u) + \text{const}$, a hot
replica's $\beta_i \ll 1$ compresses every log-likelihood difference by the same
factor: a barrier between two modes, sharp at $\beta=1$, becomes shallow enough for
an ordinary random walk to cross once $\beta_i$ is small enough --- the mechanism
behind \S\ref{sec:pt-what}'s ``flattened likelihood surface.''

Treat every replica together as a single chain on one joint space ---
\yamlkey{num\_chains} copies of the unit cube, one per rung --- with joint target
$\Pi(u_0, u_1, \ldots) = \prod_i \pi_i(u_i)$; replicas are independent under $\Pi$
because each $\pi_i$ only involves its own coordinate. Every ordinary Metropolis step
updates one replica's own coordinate using only that replica's own $\pi_i$
(Chapter~\ref{ch:toymcmc}'s rule, run once per replica each generation), so by the
previous paragraph it leaves that coordinate's $\pi_i$ invariant while leaving every
other coordinate untouched --- which leaves the whole product $\Pi$ invariant too.
The extended ensemble is where parallel tempering actually lives: $\Pi$, not any
single $\pi_i$, is the distribution the combined algorithm is built to
preserve\cite{EarlDeem:2005}.

The exchange move is a second kernel on that same joint space: it proposes to swap
coordinates $i$ and $j$ --- $(\ldots, u_i, \ldots, u_j, \ldots) \to (\ldots, u_j,
\ldots, u_i, \ldots)$ --- leaving every other coordinate fixed. Swapping is its own
inverse, so the proposal is symmetric and Chapter~\ref{ch:toymcmc}'s
symmetric-proposal simplification applies again: accept with probability
$\min(1,\ \Pi(\text{swapped})/\Pi(\text{current}))$. Every factor of $\Pi$ except
$\pi_i$ and $\pi_j$ is identical on both sides of that ratio and cancels, leaving
$\pi_i(u_j)\,\pi_j(u_i)\,/\,[\pi_i(u_i)\,\pi_j(u_j)] =
L(u_j)^{\beta_i}L(u_i)^{\beta_j}\,/\,[L(u_i)^{\beta_i}L(u_j)^{\beta_j}]$, whose log is
exactly $\Delta = (\beta_i-\beta_j)(\log L_j - \log L_i)$ ---
\S\ref{sec:pt-algorithm}'s swap rule, not asserted this time but
derived\cite{Geyer:1991,HukushimaNemoto:1996}.

Ordinary steps and exchanges are two different kernels that each leave the same
$\Pi$ invariant, so alternating between them --- exactly what
\S\ref{sec:ptmcmc-runtime}'s generation loop does --- leaves it invariant too. That
is the formal version of \S\ref{sec:pt-what}'s warning: the extended chain's
stationary distribution is the product $\Pi$ over \emph{every} rung, and only its
cold-replica marginal, $\pi_0(u) \propto L(u)$, is the posterior you asked for.

\section{How the code implements it}
\label{sec:ptmcmc-classes}

\begin{quote}
\sffamily\footnotesize
\code{MCMCBaseSampler} $\;\to\;$ \code{PTMCMCBase} $\;\to\;$ \code{PTMCMCSampler}
\end{quote}

\code{PTMCMCBase} adds exactly the parallel-tempering knobs on top of
Chapter~\ref{ch:toymcmc}'s base class: the temperature ladder, the exchange interval,
adjacent-pair selection, and the extra checkpoint state
(\S\ref{sec:ptmcmc-resume}). \code{PTMCMCSampler} itself is, once again, a thin
subclass --- it normalises \yamlkey{proposal\_scale} and attaches the progress
observer; the exchange science lives entirely in the shared base.

\subsection{The runtime: barrier-coupled generations}
\label{sec:ptmcmc-runtime}

Parallel tempering cannot use Chapter~\ref{ch:toymcmc}'s async independent-chain
pipeline (\S\ref{sec:mcmc-runtime}) --- an exchange needs every replica's
log-likelihood \emph{at the same step}, so replicas cannot run ahead of each other
independently. \sampler{PTMCMC} instead uses the classical generation-barrier path:

\begin{enumerate}
  \item Propose one step for \emph{every} replica; submit them all to the shared
        \code{hep:feedback} queue (not the per-chain shards
        Chapter~\ref{ch:toymcmc} uses).
  \item Wait for the whole generation's results; absorb every replica's
        accept/reject decision.
  \item If \yamlkey{exchange\_interval} draws have passed since the last attempt,
        attempt exchanges (\S\ref{sec:pt-algorithm}).
  \item Checkpoint at this generation barrier; repeat.
\end{enumerate}

Every generation evaluates all \yamlkey{num\_chains} replicas, even though only the
cold one is your posterior --- parallel tempering trades that extra compute for the
ability to cross barriers a single chain cannot.

\section{Configuring PTMCMC}
\label{sec:ptmcmc-yaml}

The minimal working card:

\begin{yamlwin}
Sampling:
  Method: PTMCMC
  Variables:
    - name: x
      distribution: {type: Flat, parameters: {min: 0.0, max: 1.0}}
    - name: y
      distribution: {type: Flat, parameters: {min: 0.0, max: 1.0}}
  Bounds:
    num_chains: 4
    num_iters: 2000
    proposal_scale: 0.2
    temperature_ladder: [1.0, 2.0, 4.0, 8.0]
    exchange_interval: 1
    seed: 0
  LogLikelihood:
    - {name: LogL, expression: "..."}   # your physics
\end{yamlwin}

Table~\ref{tab:ptmcmc-bounds} lists the keys on top of
Chapter~\ref{ch:toymcmc}'s \yamlkey{num\_chains} / \yamlkey{num\_iters} /
\yamlkey{proposal\_scale} / \yamlkey{seed} (unchanged in meaning here).

\begin{table}[ht]
  \footnotesize
  \begingroup
  \setlength{\tabcolsep}{0pt}
  \begin{tabular}{@{}p{0.255\linewidth}p{0.161\linewidth}p{0.584\linewidth}@{}}
    \toprule
    \textbf{Key} & \textbf{Default} & \textbf{Meaning} \\
    \midrule
    \yamlkey{temperature\_ladder} & \req & one temperature per replica
      (\S\ref{sec:pt-what}); length must equal \yamlkey{num\_chains}, strictly
      increasing, first entry exactly $1.0$ \\
    \yamlkey{exchange\_interval} & \code{1} & draws between exchange attempts
      (\S\ref{sec:pt-algorithm}) \\
    \bottomrule
  \end{tabular}
  \endgroup
  \caption{\sampler{PTMCMC}'s own \yamlkey{Bounds} keys. \yamlkey{num\_chains} here
  additionally requires $\geq 2$ (one cold, at least one hot). \code{Jarvis man
  sampler.PTMCMC} is the live reference.}
  \label{tab:ptmcmc-bounds}
\end{table}

\begin{jtip}
A common starting ladder is geometric --- $T_i = r^{\,i}$ for some ratio $r$ between
$1.5$ and $2$ --- which is exactly what \yamlkey{[1.0, 2.0, 4.0, 8.0]} above is. Watch
the exchange acceptance rate once a run is going (\S\ref{sec:ptmcmc-log-progress}): if
adjacent rungs almost never swap, they are too far apart on the ladder; if they almost
always swap, the rungs are barely distinguishable and you are paying for replicas that
add little.
\end{jtip}

\section{Validating your card}
\label{sec:ptmcmc-validation}

\begin{terminal}
\begin{jout}
Config validation failed (2 error(s), 0 warning(s)):

Error summary (reference only):
+---+-------------+------------------------------------+-------------------------+
| # | Code        | YAML path                          | Problem                 |
+---+-------------+------------------------------------+-------------------------+
| 1 | JV2-SCH-001 | \$.Sampling.Bounds                  | 'temperature_ladder'... |
| 2 | JV2-MTH-064 | Sampling.Bounds.temperature_ladder | PTMCMC requires Sam...  |
+---+-------------+------------------------------------+-------------------------+
This table is only a quick reference. Please consult the Jarvis-HEP YAML settings
documentation before editing the task card.

Detailed diagnostics:
  [error] JV2-SCH-001  \$.Sampling.Bounds
          'temperature_ladder' is a required property
          hint: Run: Jarvis man yaml.Sampling.Bounds
          suggestion: Add 'temperature_ladder' at this location using the
          expected YAML type.
  [error] JV2-MTH-064  Sampling.Bounds.temperature_ladder
          PTMCMC requires Sampling.Bounds.temperature_ladder (a strictly
          increasing list starting at 1.0)
          hint: Run: Jarvis man sampler.PTMCMC
          suggestion: Set temperature_ladder to a strictly increasing list of
          length num_chains starting at 1.0.
          example:
            Bounds:
              num_chains: 4
              temperature_ladder: [1.0, 2.0, 4.0, 8.0]
\end{jout}
\end{terminal}

Table~\ref{tab:ptmcmc-validation} lists the \sampler{PTMCMC}-specific codes, on top of
Chapter~\ref{ch:toymcmc}'s \code{JV2-MTH-05x} family for
\yamlkey{num\_chains}/\yamlkey{num\_iters}/\yamlkey{proposal\_scale}.

\begin{table}[ht]
  \footnotesize
  \begingroup
  \setlength{\tabcolsep}{0pt}
  \begin{tabular}{@{}p{0.171\linewidth}p{0.297\linewidth}p{0.532\linewidth}@{}}
    \toprule
    \textbf{Code} & \textbf{Fires on} & \textbf{Fix} \\
    \midrule
    \code{JV2-MTH-060} & \yamlkey{Bounds} missing entirely & add all four required
      keys \\
    \code{JV2-MTH-061} & \yamlkey{num\_chains} missing or $< 2$ & \sampler{PTMCMC}
      needs at least one hot replica \\
    \code{JV2-MTH-064} & \yamlkey{temperature\_ladder} missing, wrong length, not
      increasing, or not starting at $1.0$ & fix the ladder
      (\S\ref{sec:ptmcmc-yaml}) \\
    \code{JV2-MTH-066}--\code{071} & other malformed keys (\yamlkey{exchange\_interval},
      \yamlkey{seed}, \ldots) & check against Table~\ref{tab:ptmcmc-bounds} \\
    \bottomrule
  \end{tabular}
  \endgroup
  \caption{\sampler{PTMCMC}-specific validation codes. The other methods' codes are
  in Appendix~\ref{app:troubleshooting}.}
  \label{tab:ptmcmc-validation}
\end{table}

\section{Reading the log}
\label{sec:ptmcmc-log}

Every excerpt below is a real \cli{Jarvis run} transcript against the same
two-parameter Gaussian bump as Chapter~\ref{ch:toymcmc}, with
\yamlkey{temperature\_ladder: [1.0, 2.0, 4.0, 8.0]}.

\subsection{At startup}
\label{sec:ptmcmc-log-startup}

\begin{terminal}[sampler.log]
\begin{jout}
·•· Jarvis-HEP.Sampler.Ptmcmc
	-> 00:15:16.099 - [WARNING] >>>
Initializing the PTMCMC Sampling

·•· Jarvis-HEP.Sampler.Ptmcmc
	-> 00:15:16.099 - [INFO] >>>
PTMCMC Sampler configured: replicas=4 iterations=50 total_transitions=200
proposal_scale=0.2 ladder=[1.0, 2.0, 4.0, 8.0] exchange_interval=1
\end{jout}
\end{terminal}

Confirm \yamlkey{ladder} here matches your card exactly --- this is the resolved
configuration actually driving the run, the same discipline as nested sampling's
\code{run\_nested kwargs} line (Section~\ref{sec:nested-log-startup}).

\subsection{While it runs: exchanges and the progress meter}
\label{sec:ptmcmc-log-progress}

Exchange attempts log immediately, before the progress line for that generation:

\begin{terminal}[sampler.log]
\begin{jout}
·•· Jarvis-HEP.Sampler.Ptmcmc
	-> 00:15:16.542 - [INFO] >>>
PTMCMC exchange: attempted=2 accepted=1

·•· Jarvis-HEP.Sampler.Ptmcmc
	-> 00:15:16.542 - [WARNING] >>>
20‰ of 4/200 PTMCMC transitions completed in 00:00:00.442
(replicas=4 iterations=1/50 swaps=1/2)
\end{jout}
\end{terminal}

\code{swaps=1/2} in the progress line is the running total (accepted/attempted) since
the run started --- by the end of this $200$-transition run it read
\code{swaps=49/75}, a $65\,\%$ exchange acceptance rate across the whole ladder.

\subsection{At the end: the summary}
\label{sec:ptmcmc-log-summary}

\begin{terminal}[sampler.log]
\begin{jout}
·•· Jarvis-HEP.Sampler.Ptmcmc
	-> 00:15:16.970 - [INFO] >>>
PTMCMC finished: proposed=200 accepted=121 accept_rate=0.6050 async=False
\end{jout}
\end{terminal}

\code{async=False} is \Jtwo\ confirming this run used the barrier-coupled pipeline
(\S\ref{sec:ptmcmc-runtime}), not Chapter~\ref{ch:toymcmc}'s async one. The overall
\yamlkey{accept\_rate} here counts ordinary Metropolis steps across all four
replicas --- hot replicas accept far more often than the cold one by design (this run:
$30\,\%$, $48\,\%$, $78\,\%$, $86\,\%$ from coldest to hottest), which is exactly the
flattened-surface behaviour \S\ref{sec:pt-what} describes, not a sign anything is
wrong.

\section{What the run leaves behind}
\label{sec:ptmcmc-outputs}

The same two files as Chapter~\ref{ch:toymcmc} (\S\ref{sec:mcmc-outputs}), with three
extra fields in \file{DATABASE/sampler\_summary.json}:

\begin{table}[ht]
  \footnotesize
  \begingroup
  \setlength{\tabcolsep}{0pt}
  \begin{tabular}{@{}p{0.255\linewidth}p{0.745\linewidth}@{}}
    \toprule
    \textbf{Field} & \textbf{Meaning} \\
    \midrule
    \yamlkey{swap\_attempts} / \yamlkey{swap\_accepts} & running totals across the
      whole run --- their ratio is the overall exchange acceptance rate
      (\S\ref{sec:ptmcmc-yaml}'s tuning tip) \\
    \yamlkey{temperature\_ladder} & the resolved ladder actually used, for the
      record \\
    \bottomrule
  \end{tabular}
  \endgroup
  \caption{\sampler{PTMCMC}-only \code{sampler\_summary.json} fields.}
  \label{tab:ptmcmc-summary-fields}
\end{table}

\yamlkey{chain\_id}'s \yamlkey{temperature} in \file{chain\_history.csv} (and each
chain's entry in \yamlkey{sampler\_summary.json}'s \yamlkey{chains} list) is fixed to
that \yamlkey{chain\_id}'s ladder slot for the whole run --- it is the
\emph{configuration} that migrates between rungs on a swap, never the temperature
label itself (\S\ref{sec:pt-algorithm}).

\section{Resume}
\label{sec:ptmcmc-resume}

Everything in Chapter~\ref{ch:toymcmc}'s \S\ref{sec:mcmc-resume} applies --- the same
\file{state.pkl}, the same \cli{Jarvis run scan.yaml --resume}, the same refusal on a
changed card. \sampler{PTMCMC} checkpoints one generation barrier at a time (never
mid-generation, unlike the async pipeline) and additionally restores
\yamlkey{swap\_attempts}, \yamlkey{swap\_accepts}, the exchange pairing offset, and the
resolved \yamlkey{temperature\_ladder} --- so a resumed run's exchange statistics keep
accumulating from where the interrupted one left off, rather than restarting at zero.

\section{Summary}
\label{sec:ptmcmc-summary}

\sampler{PTMCMC} is Chapter~\ref{ch:toymcmc}'s Metropolis chain, run at several
temperatures with periodic exchanges --- the answer when a posterior might be
multimodal enough to trap a single chain. Table~\ref{tab:ptmcmc-at-a-glance} gathers
it.

\begin{table}[ht]
  \footnotesize
  \begingroup
  \setlength{\tabcolsep}{0pt}
  \begin{tabular}{@{}p{0.34\linewidth}p{0.66\linewidth}@{}}
    \toprule
    \textbf{Field} & \textbf{\sampler{PTMCMC}} \\
    \midrule
    Kind & feedback (fixed budget: replicas $\times$ iterations) \\
    Pipeline & barrier-coupled generations (\S\ref{sec:ptmcmc-runtime}) \\
    Point count & exactly \yamlkey{num\_chains}$\,\times\,$\yamlkey{num\_iters}, only
      \yamlkey{chain\_id 0} is posterior-valid \\
    Dimensions & any \\
    Required keys & \yamlkey{num\_chains} ($\geq 2$), \yamlkey{num\_iters},
      \yamlkey{proposal\_scale}, \yamlkey{temperature\_ladder} \\
    u$\to$x mapping & yes \\
    \yamlkey{Bounds.seed} role & chain RNG streams and the exchange \textsc{rng} \\
    \textsc{uuid}s & deterministic, as Chapter~\ref{ch:toymcmc} \\
    Mid-run resume & yes --- one generation barrier at a time \\
    Extra output & \file{sampler\_summary.json} (+ swap stats),
      \file{chain\_history.csv} \\
    \bottomrule
  \end{tabular}
  \endgroup
  \caption{\sampler{PTMCMC} at a glance.}
  \label{tab:ptmcmc-at-a-glance}
\end{table}
